\section{一张清丈册能看见什么}
\label{sec:wanli-survey}

万历初年的朝廷试图把缓慢的地方政务放进可以催问的期限。考成法让部院和地方官必须回复“办到哪里”，却不等于北京从此知道事情在每一个县怎样发生。报送及时可能意味着工程、征解和裁决确有进展，也可能意味着官员更熟练地安排文书。张居正的改革力量，正在于它把延误、空额和互相推诿变成可以追问的行政问题；它的边界也正在于账册仍要由地方书吏、业主、里甲与州县共同写成。

\eventref{MM-1578-NATIONWIDE-SURVEY}{清丈}把这个问题推到田亩上。中央的命令、各省的启动和报亩，可以由实录、户部材料和地方赋役册互校；不同地区的完成年月、丈量方法和隐漏争议却不能压成同一年、同一尺度的“全国实测”。一张田册把土地变成应比较、应征收的对象，也把产权、佃作、逃役和地方中介的冲突写得更尖锐。它不是对乡村的摄影，而是一项为了取得财政可见性而进行的行政劳动。

\eventref{MM-1580-ONE-WHIP-ORDERS}{一条鞭的推行}紧接着改变若干田赋、丁役和杂派的计征与折纳方式。这里的“银”不是抽象的富裕，也不是所有人手中同样可得的货币。江南市镇、矿区、商路和海外贸易让银的流动更密集；边镇、灾区和以实物、差役为重的地区，则可能以不同的兑换、借贷和摊派承受政策。更多责任转向田亩，并不表示人丁、劳役、实物和临时加派从社会中消失。改革把征收链条重新打结，不能消除链条上每一环的谈判与损耗。

\section{三处战场与一条粮道}
\label{sec:wanli-three-campaigns}

万历二十年，宁夏的军事冲突与日本入侵朝鲜几乎同时要求兵员、粮饷和判断。对朝廷而言，调兵清单可以排列得很整齐；对每一镇而言，出发意味着谁留下守堡、粮从哪里解、船和道路能否承载、欠饷如何补上。明朝介入朝鲜的决定不能只写成“援朝”二字。朝鲜王朝需要保存自身政权与民众，日本各军事集团有不同的补给和扩张目标，明军也不只从北京的一纸诏令获得行动能力。

《神宗实录》、《朝鲜王朝实录》、明清正史和日本、海上材料各自给出战事的时间骨架，却对兵数、战果、后勤和责任说法不一。它们的差异不是阅读的障碍，而是问题本身：战报要争取资源，失败后的叙事要分配责任，地方记录则未必看见高层的全部计划。战争最终以日军撤离收束，并不意味着军费已经回到原有水平。朝鲜的毁坏、边镇的调动、东南海防和中央的债欠，使“胜利”带着长久的财政余震。

因此，万历后期的军费不能由一个崩溃数字解释。宁夏、朝鲜与随后辽东的花费来自不同的粮道、银库、临时加派和地方摊解。法定额、请拨额、报解额与实际抵达的资源必须分列；任何一份总账都同时是行政愿望和文书策略。正是这种差距，让朝廷在新的战争到来时越来越难以把一条命令变成同时到达多处前线的粮饷。

\section{宫中、书院与乡里的政治时间}
\label{sec:wanli-court-local}

张居正去世以后，清丈与一条鞭没有随同一个人立即消失；它们已被写进许多地方的征收实践，同时失去了能协调部院、皇帝和地方官的强力中枢。万历朝后期的停朝、任官争执和国本问题，常被叙述成远离民生的宫廷戏。其实，部院不能正常补缺、奏疏不能及时裁决、税饷与河工的责任无人协调，会让县以下的等待、借贷和摊派有不同的结果。宫门里的政治时间，经过文书与人事，仍会抵达运河闸口、仓库和乡村。

书院限制也显示国家与地方知识网络的拉扯。禁令能说明中央对讲学、结社和声望聚集的担忧；碑记、文集和地方志才能较有限地追踪一处书院是停办、改名、转入家塾，还是在另一个名义下继续。士人并不是只在朝廷与乡里两端择一，他们借助宗族、同年、学校、商人资助与地方公议连接两端。将这些网络统称为“党争”，既过早替人规定身份，也难以解释地方公共事务为何能在不同处所持续运作。

\section{辽东不是一条突然断裂的边线}
\label{sec:wanli-liaodong}

后金的建立、抚顺失守和萨尔浒战败，使辽东的问题以城堡和战报的形式压到朝廷。城池失守可以确认，谁掌握了市场、卫所、女真部众与朝鲜方向的关系，却不能由一方军报写成单线故事。努尔哈赤的政权并非从一开始就是内部无差别的整体；明朝辽东也不是一堵只有军队的墙。卫所家庭、军商、屯田、翻译和各族群关系，使防线既有军事面，也有社会与交换的面。

萨尔浒之后，失去的不只是几支部队，更是协调多路进军、粮道和情报的信心。兵数、伤亡和路线在明方记录、后出正史与对方叙事中差异很大，应当保留，而非选择一个最戏剧性的版本。万历末年的连续继承又令处理辽东的时间更加压缩。下一章要面对的，正是已经被军费、灾害、道路和社会流动同时拉紧的帝国，而不是一张可以由“白银不足”或“党争激烈”任意解释的简图。
